Praktika honen helburu nagusia mugikorreko roboten kontrolerako oinarrizkoa den funtzio anitzen arrastoa eta atzeraelikadura lantzea da. Oraingoan, paretari jarraitzea aztertuko dugu laser bidezko erregresio linealarekin eta PID kontrolagailuarekin. Robota eskuineko paretari paraleloki jarraituz higituko da, erregresio bidez kalkulatutako $\theta$ orientazio-angelua $0^\circ$-ra eramanez.

\needspace{6\baselineskip}
\section{Testuingurua eta motibazioa}\label{sec:testuingurua}
Robotika mugikor autonomoan, inguruneko egitura finkoak (hala nola paretak, korridoreak edo mugak) erreferentzia gisa hartzea ezinbesteko estrategia da ibilbide egonkorrak osatzeko. Aurreko praktikan kameraren bidezko bide-jarraipen bisuala aztertu genuen bezala, oraingoan 2D LiDAR sentsoreak ematen duen distantzia-informazio geometrikoa erabiliko dugu. Erronka nagusia da sentsorearen neurketa gordinetatik abiatuta paretaren orientazio erlatiboa modu sendoan estimatzea eta informazio horrekin atzeraelikadurazko kontrol-begizta bat ixtea, talkarik edo oszilaziorik gabe.

\needspace{6\baselineskip}
\section{Kontrol-estrategia eta printzipioak}\label{sec:estrategia}
Praktika honetan \textbf{HIPER-KISS} metodologia zorrotza jarraitu dugu: irakasleak eskatutako oinarrizko baliabideak soilik erabili dira, alferrikako abstrakziorik eta mekanismo gehigarririk gabe:
\begin{itemize}[leftmargin=2em, itemsep=0.3em]
  \item \textbf{Errore garbia}: Kontrolatzailearen errorea $\theta$ angelua da soilik ($e = \theta - 0{,}0$), helburu bakarra paretarekiko paralelotasuna baita. Ez da distantzia-termino artifizialik sartu kontrol-legean.
  \item \textbf{Segurtasun-arau frontalik gabe}: Ez da aurrealdeko balaztatze- edo biraketa-salbuespenik programatu; izkinak PID kontrolagailuak berak kudeatzen ditu hautatutako izpi-sektoreari esker.
  \item \textbf{Izpien sektorea}: Eskuin-aurreko koadrantea hartu da ($-90^\circ$-tik $0^\circ$-ra). Sektorea aurrerantz $0^\circ$-raino iristen denez, aurreko pareta $d_{\max}$ distantziara iristean sartzen da erregresioan, puntuek ``L'' forma hartzen dute, malda positiboa sortzen da eta robotak berez ezkerrera biratzen du izkinetan.
  \item \textbf{$d_{\max}$ muga}: Irakurketa motzen mugak erabakitzen du zein distantziatan jarraituko dion robotak paretari. Izkina bakoitzean $d_{\max}$ inguruan hasten du biraketa, eta distantzia $D \approx d_{\max} - 0{,}25\text{ m}$-tan egonkortzen da.
  \item \textbf{Abiadura lineala}: $v = v_0 / (1 + k_v \cdot |\omega|)$ formula aplikatu da. $k_v = 0$ denean abiadura konstantea da; $k_v > 0$ denean, biraketetan moteltzen da.
\end{itemize}

\needspace{6\baselineskip}
\section{Emaitzen jatorria}\label{sec:emaitzen-jatorria}
Txosten honetako taula eta grafiko guztiak ROS 2 Humble nodo errealarekin eta proba-simulagailu zinematiko batekin neurtu dira ($5 \times 3{,}5\text{ m}$-ko pareta laukizuzena, $T = 0{,}1\text{ s}$ laginketa-periodoa). Emaitza hauek azterketa kualitatibo eta kuantitatibo sendoa ahalbidetzen dute, inertzia eta marruskadura konplexurik gabeko erreferentzia garbi gisa. Gazeboko mundu birtualean neurtzean, datuak \texttt{wf\_plot.py} script-aren bidez birsortu ahal izango dira.

\needspace{6\baselineskip}
\section{Dokumentuaren egitura}\label{sec:egitura}
Txostena ondorengo ataletan banatzen da:
\begin{itemize}[leftmargin=2em, itemsep=0.3em]
  \item \ref{ch:garapena}. kapituluan LiDAR neurketen proiekzio kartesiarra, erregresio lineala, P kontrolaren dinamika diskretua eta PID kontrolagailu osoaren formulazio matematikoa aurkezten dira.
  \item \ref{ch:emaitzak}. kapituluan parametro-barridoen emaitzak ($K_p$, $K_i$, $K_d$, abiadura-legeak), ibilbide-grafikak eta hiru aldaera nagusien arteko konparaketa zehatza azaltzen dira.
  \item \ref{ch:ondorioak}. kapituluan lortutako ondorio nagusiak eta Gazeborako gomendio praktikoak laburbiltzen dira.
  \item \ref{ch:eranskina}. eranskinean simulazio-aginduak eta garatutako fitxategien katalogoa jasotzen dira.
\end{itemize}
